\documentclass[11pt,letterpaper]{article}
\usepackage[margin=0.60in,top=0.58in,bottom=0.56in]{geometry}
\usepackage{mathptmx}
\usepackage{microtype}
\usepackage{graphicx}
\usepackage{booktabs}
\usepackage{array}
\usepackage{longtable}
\usepackage{amsmath,amssymb}
\usepackage[hidelinks]{hyperref}
\setlength{\parindent}{0pt}
\setlength{\parskip}{2.5pt}
\setlength{\columnsep}{0.25in}
\setcounter{secnumdepth}{0}
\renewcommand{\section}[1]{\vspace{3pt}\noindent\textbf{\underline{#1}}\par\vspace{1pt}}
\renewcommand{\subsection}[1]{\vspace{2pt}\noindent\textbf{#1}\par\vspace{1pt}}
\renewcommand{\familydefault}{\rmdefault}
\providecommand{\tightlist}{\setlength{\itemsep}{0pt}\setlength{\parskip}{0pt}}
\providecommand{\pandocbounded}[1]{\sbox0{#1}\ifdim\wd0>\linewidth\resizebox{\linewidth}{!}{#1}\else\usebox0\fi}
\title{\vspace{-1.2em}\textbf{\MakeUppercase{Neurologically Motivated
Self-Supervised Learning for Bilateral Gait Geometry}}\vspace{-0.5em}}
\author{\normalsize [Authors and affiliations to be inserted in the RISEx template]}
\date{}
\begin{document}
\twocolumn[
  \begin{@twocolumnfalse}
  \maketitle
  \vspace{-1.1em}
  \end{@twocolumnfalse}
]
\subsection{INTRODUCTION}\label{introduction}

Self-supervised motion features are attractive for AI systems that must
interpret people from incomplete observations. A central question is
whether success at predicting a hidden feature means that the
representation retains a useful physical relation. We test this question
with bilateral gait geometry. Reflection exchanges left and right body
landmarks and reverses a signed contrast of their movement speeds. This
relation is relevant to gait asymmetry in several neurological and
musculoskeletal conditions {[}1{]}, while remaining independent of a
diagnostic label.

\subsection{MATERIALS AND METHODS}\label{materials-and-methods}

The study uses 625 quality-checked GAVD pose clips from 93 source videos
{[}2{]}. Each clip is prepared independently. All 33 landmarks remain in
a 64-by-33-by-3 input, and four-frame patches yield a 16-by-33
time-by-joint grid. A skeleton JEPA predicts teacher features for masked
grid positions from their visible context {[}3{]}. No label, target, or
affected side supervises encoder training.

We define a left-minus-right median-speed contrast over five paired
joints. The null hypothesis is that JEPA does not improve its frozen
ridge-regression estimate over a matched initial encoder. Five outer
folds exclude complete source videos from pretraining and readout
fitting; five seeds repeat each condition. Mask comparisons share
starting weights, source draws, training views, and updates. We weight
sources equally and resample sources, not clips, for paired intervals.

\subsection{RESULTS AND DISCUSSION}\label{results-and-discussion}

{\def\LTcaptype{none} % do not increment counter
\begin{tabular}{@{}p{0.40\columnwidth}p{0.22\columnwidth}p{0.28\columnwidth}@{}}
\toprule\noalign{}
Outcome & Before JEPA & After JEPA \\
\midrule\noalign{}
Bilateral readout \(R^2\) & 0.223 & 0.101--0.114 \\
Correct-clip target preferred & 33/75 checks & 375/375 checks \\
\end{tabular}
}

All five trained-minus-initial contrasts were negative (\(-0.109\) to
\(-0.122\)); their saved source intervals lay below zero.
Motion-weighted and connected-region masks had no clear advantage over
their own count-matched random references. The predictor therefore
learns within-clip feature correspondence while the tested
encoder-plus-readout pipeline produces a weaker bilateral movement
estimate.

The evidence has limits that matter for interpretation. Teacher feature
scales differ between conditions. The expanded readout adds temporal
summaries and missing-data support while increasing dimension. Moreover,
preparation changes the target calculation: its agreement with the
original calculation is 70.4\% by sign and \(R^2=0.218\). The study does
not show that JEPA destroys geometric information, only that this
training recipe and readout fail to improve this observable.

\subsection{CONCLUSIONS}\label{conclusions}

A latent prediction diagnostic can improve even when a prespecified
geometric readout declines. Human-aware robotics may benefit from
testing learned motion features against task-relevant observables before
using them downstream. Future work should evaluate past-only motion
prediction with persistence and velocity baselines, a stable target
path, and new source or participant groups.

\subsection{REFERENCES}\label{references}

{[}1{]} Lewek MD et al.~\emph{Gait Posture} 31:256--260, 2010. {[}2{]}
Ranjan R et al.~\emph{IEEE Access} 13:45321--45339, 2025. {[}3{]}
Abdelfattah M, Alahi A. \emph{ECCV}, 2024.
\end{document}
